\scanpage{134}

1\deleted{2}\handwritten{3}) Variante utile du tapis de n° 6. Soit $C$ une catégorie, et $\deleted{\unsure{F:C\;\mathrm{Mor}|F:C\;\mathrm{Mo}|F:C\;\mathrm{Mod}}}$
\[
F:C\handwritten{\longrightarrow}\mathrm{Modf}(k)
\]
un foncteur, où $\deleted{j}k$ est un corps. On ne suppose pas $C$ additive, ni $F$ fidèle. On définit alors l'anneau profini $U=\mathrm{End}_{k}(F)$, la coalgèbre duale $A$, et on trouve une factorisation canonique de $F$ à travers $\mathrm{Modf}(U/k)=\mathrm{Comodf}(A/k)$. Supposons donné dans $C$ un bifoncteur $\otimes$, \deleted{a\unsure{\deleted{t}l|l}ors}\inserted{\handwritten{et}} un isomorphisme de bifoncteurs $F(X\otimes Y)\handwritten{\simeq}F(X)\otimes F(Y)$, alors on trouve une application diagonale sur $U$ i.e. une loi de composition \deleted{\unsure{comp|comp.} avec \unsure{appli|appl.}} sur $A$: en effet, on constate que l'anneau des endomorphismes du bifoncteur du premier membre est $U\otimes_{k}U$, et $U$ s\deleted{e}'envoie dedans grâce à l'isomorphisme donné $F(X\otimes Y)\handwritten{\simeq}F(X)\otimes F(Y)$. Si on se donne pour $\otimes$ des données d'unité, d'associativité ou de commutativité, la compatibilité de ces données avec $c$ s'exprime alors comme il est dit dans n° 6\inserted{\handwritten{en \unsure{termes|terme} de $\Delta$}}\figure{\begin{tikzpicture}[baseline=-.6ex,x=1cm,y=1cm]
\draw (.06,-.24) .. controls (.02,.12) and (.34,.18) .. (.75,.2) -- (4.1,.2) .. controls (4.5,.22) and (4.5,.32) .. (4.46,.43);
\end{tikzpicture}}, et il y a lieu de donner des variantes graduées (si possible plus générale\handwritten{s} et astucieuse\handwritten{s} qu'au n° 8\handwritten{!)}. En fait, les résultats spécifiques au n°6 seraient \unsure{pl\deleted{û}utôt|pl\deleted{t}ûtôt|plûtôt}, dans le cas d'un foncteur \unsure{fi\deleted{è}d\deleted{d}èle|fi\deleted{d}è\deleted{è}dle|fi\deleted{d}è\deleted{d}èle} (et additif \deleted{de c\unsure{a|o}té} exact de catégories abéliennes ? à voir de quoi exactement on a besoin), des réciproques aux constructions précédentes, i.e. comment déduire d'une application diagonale sur $U$ un\deleted{e} foncteur $\otimes$ sur $C$, do\deleted{j}nt les pr\deleted{i}opriétés vont réfléter celles de cette dernière.

Exemple intéressant: les opérations cohomologiques sur $\underline{F}_{p}$. Il mérite quelques explications: sauf le morceau engendré par les $\handwritten{\pi}_{i}$ \deleted{(NB $i$} les degrés dans $U$ sont $\handwritten{>}$ $0$, et $\overset{\handwritten{\sim}}{U}\inserted{\handwritten{=k+U^{+}}}\text{\figure{\begin{tikzpicture}[baseline=-.6ex,x=1cm,y=1cm]
\draw (.06,-.24) .. controls (.02,.12) and (.34,.18) .. (.75,.2) -- (2.4,.2) .. controls (2.8,.22) and (2.8,.32) .. (2.76,.43);
\end{tikzpicture}}}$ est le complété c\unsure{\deleted{n}a\deleted{a}n|\deleted{n}a\deleted{n}a|\deleted{a}n\deleted{n}a}onique d'un anneau gradué à application diagonale anticommutative, donc $\overset{\handwritten{\sim}}{A}$ est un anneau analogue. Dire comment on récupère $U,A$ à partir de $\overset{\handwritten{\sim}}{U}$, $\overset{\handwritten{\unsure{\sim}}}{\overset{\handwritten{\sim}}{A}}$.
